\section{Monomial spaces and finite jets}
\label{sec:transfers}

The proof will replace spaces of sections by monomial spaces and then change
their exponent sets. Two properties must survive these changes: multiplication
between different degrees, and the implication that interpolation for the new
space yields interpolation for the old one. We establish these properties
separately. Multiplication concerns the exact exponent sets in every degree;
interpolation will be transferred by analytic limits at one fixed degree.

\subsection{Jets and graded exponent sets}

Write $\Torus^n=(\C^*)^n$. For a finite set $S\subset\Z^n$, put
\[
 \M(S)=\Span_{\C}\{x^\alpha:\alpha\in S\},
 \qquad x^\alpha=x_1^{\alpha_1}\cdots x_n^{\alpha_n}.
\]
These Laurent polynomials are regular on $\Torus^n$. If $A$ is a line bundle
on a smooth complex variety $Y$, its jets through an integer degree $m\geq0$ at $p\in Y$
form the vector space
\[
 \Jet_p^m(A)
 =A_p/\frakm_p^{m+1}A_p.
\]
Here $A_p$ denotes the stalk of the sheaf of sections. A space of sections
$W$ \emph{separates jets through degree $m$} at distinct points
$p_1,\ldots,p_r$ if the evaluation map
\[
 W\longrightarrow\bigoplus_{i=1}^r\Jet_{p_i}^m(A)
\]
is surjective. For functions we use the trivial line bundle. Local frames and
local coordinates change the matrices of these maps but not their ranks.
Analytic and algebraic finite jets agree at a smooth complex point, since
the corresponding local rings have the same completion.
Our aim is to obtain such surjections for $W=H^0(X,kL)$ with $m/k$
approaching $(V/r)^{1/n}$.

\begin{definition}\label{def:graded-support}
A \emph{graded support family} in $\Z^n$ is a sequence of finite sets
$(S_k)_{k\geq1}$ such that
\begin{equation}\label{eq:graded-support}
 S_k+S_l\subset S_{k+l}\qquad(k,l\geq1).
\end{equation}
Empty sets are allowed. Equivalently,
$\M(S_k)\M(S_l)\subset\M(S_{k+l})$, where a product of spaces denotes the
span of their pairwise products.
\end{definition}

We will say that an operation on monomial spaces \emph{transfers jet
separation backward} if, for every fixed degree, number of points, and jet
order, separation by the output space at some distinct torus tuple implies
separation by the input space at some distinct torus tuple. The tuples need
not be the same. A finite sequence of such operations has the same property.

\subsection{Least terms and analytic limits}

A monomial order on $\Z_{\geq0}^q$ will mean a total well-order preserved by
addition. For a nonzero convergent power series $f$, let $\nu(f)$ be the
least exponent in its support for this order.

\begin{lemma}[Least-term bases]\label{lem:least-terms}
Let $W$ be a finite-dimensional space of convergent power series in $q$
variables, and fix a monomial order. The set
\[
 \nu(W\setminus\{0\})
\]
has cardinality $\dim W$. There is a basis of $W$ whose least exponents are
exactly this set, with each least coefficient equal to $1$. Moreover,
$\nu(fg)=\nu(f)+\nu(g)$ for nonzero convergent power series $f,g$.
\end{lemma}

\begin{proof}
If $W\ne0$, choose the smallest exponent occurring as the least exponent of
an element of $W$, and normalize one such element to have coefficient $1$
there. Continue in the kernel of that coefficient functional, which has
codimension one. After $\dim W$ steps this gives a basis with strictly
increasing least exponents. In any nonzero linear combination, the first
basis element with a nonzero coefficient determines the least term. This
proves the assertion about the whole set of possible exponents.

For the product assertion, write $a=\nu(f)$ and $b=\nu(g)$. Every exponent
of $f$ is at least $a$, and every exponent of $g$ is at least $b$. Additivity
of the order implies that any pair other than $(a,b)$ has sum strictly
greater than $a+b$. The coefficient of $x^{a+b}$ in the product is therefore
the product of the two nonzero least coefficients.
\end{proof}

\begin{lemma}[Least terms in labeled spaces]\label{lem:row-valuation}
Fix integers $1\leq q\leq n$ and a monomial order in $q$ variables. For
$k\geq1$ and $\gamma\in\Z^{n-q}$, let $W_{k,\gamma}$ be
finite-dimensional spaces of convergent power series in those variables,
with only finitely many nonzero spaces in each degree. Assume that
\[
 W_{k,\gamma}W_{l,\eta}\subset W_{k+l,\gamma+\eta}
 \qquad(k,l\geq1;\ \gamma,\eta\in\Z^{n-q}).
\]
Then
\[
 S'_k=\{(\nu(f),\gamma):0\ne f\in W_{k,\gamma}\}
\]
is a graded support family and
$\#S'_k=\sum_\gamma\dim W_{k,\gamma}$.
\end{lemma}

\begin{proof}
The cardinality assertion follows from Lemma~\ref{lem:least-terms} in each
label. Given two exponents in $S'_k$ and $S'_l$, choose functions witnessing
them. Their nonzero product belongs to the summed degree and label, and
its least exponent is the sum of the two least exponents.
Expanding products of arbitrary linear combinations therefore gives
$\M(S'_k)\M(S'_l)\subset\M(S'_{k+l})$. The full new monomial spaces
thus retain the multiplication property needed for the next operation;
no simultaneous choice of lifts to earlier spaces is required.
\end{proof}

The finite-degree passage to initial monomials is a standard
jet-degeneration principle; see \citet[Proposition~5.7]{Ito2013}. The next
two lemmas give the direct analytic form used here, including unchanged
coordinate factors. The first is analytic; the second is finite-dimensional
linear algebra.

\begin{lemma}[Weighted limits with unchanged variables]
\label{lem:weighted-limit}
Let $1\leq q\leq n$, and let $f_1,\ldots,f_N$ be convergent power series
near $0\in\C^q$. Suppose
that a vector $\lambda\in\R_{>0}^q$ makes a specified exponent $e_j$ the
unique minimum of $\lambda\cdot\alpha$ on the support of $f_j$, and let
$c_j\ne0$ be its coefficient. For any
$\gamma_j\in\Z^{n-q}$, as $s\downarrow0$ the functions
\begin{equation}\label{eq:weighted-limit}
 c_j^{-1}s^{-\lambda\cdot e_j}
 f_j(s^{\lambda_1}Z_1,\ldots,s^{\lambda_q}Z_q)U^{\gamma_j}
 \longrightarrow Z^{e_j}U^{\gamma_j}
\end{equation}
converge locally uniformly on $\C^q\times\Torus^{n-q}$, with all derivatives
of any fixed finite order. The assertion means that each fixed compact set
is in the domain of the left-hand side for all sufficiently small $s>0$.
When $q=n$, the unchanged variable block is omitted.
\end{lemma}

\begin{proof}
First omit the unchanged factor $U^{\gamma_j}$. On a bounded polydisc in
$Z$, choose $s_0>0$ small enough that the power series converges absolutely
on a slightly larger polydisc after scaling by $s_0^\lambda$. In the
normalized expansion every exponent $\lambda\cdot(\alpha-e_j)$ is
nonnegative, and every exponent other than the least one is strictly
positive. For $0<s\leq s_0$, the terms are
dominated by the summable series obtained at $s_0$. Dominated convergence
gives locally uniform convergence to $Z^{e_j}$. Applying Cauchy's estimates
on slightly larger polydiscs gives convergence of derivatives. The unchanged
Laurent monomials and their derivatives are holomorphic near every compact
subset of the torus, so multiplying by them preserves these conclusions.
\end{proof}

\begin{lemma}[Persistence of jet separation]\label{lem:rank-transfer}
Fix a finite tuple of distinct points in a complex manifold and a jet order
$m$. Suppose that finitely many holomorphic functions, depending on a
positive parameter $s$, converge with their derivatives through order $m$
near that tuple. If the limiting functions separate those jets, then the
functions for sufficiently small $s>0$ do also.

The same conclusion applies to sections written in local frames. If the
functions before the limit are pullbacks under a coordinate change sending
the tuple to distinct points, their original functions separate jets at the
image tuple. Multiplication of individual functions by nonzero constants
does not affect this statement.
\end{lemma}

\begin{proof}
Choose coordinates and frames near the points. The entries of the jet matrix
are the finitely many required Taylor coefficients. A maximal minor that is
nonzero in the limit remains nonzero for small $s$. Coordinate changes and
frame changes induce isomorphisms of truncated local rings or their
rank-one modules. Rescaling an individual function rescales a matrix column
by a nonzero constant.
\end{proof}

In particular, the unchanged variables in
Lemma~\ref{lem:weighted-limit} cause no restriction on the jets being
separated: all derivatives, including derivatives in those variables, enter
the same convergent jet matrix. We shall apply the lemmas to finitely many
spaces distinguished by a common exponent in the unchanged variables.

\subsection{Lattice changes and coordinate compression}

We next record two operations directly on Laurent monomials. They will let
us change the shape of an exponent set while retaining both properties
above.

\begin{lemma}[Unimodular changes]\label{lem:unimodular}
Let $B\in\operatorname{GL}_n(\Z)$. Replacing $S_k$ by $BS_k$ preserves
cardinality and the graded support property. It preserves, in both
directions, the existence of a distinct torus tuple at which the associated
monomial space separates any prescribed finite jets.
\end{lemma}

\begin{proof}
The lattice map is bijective and additive. The torus automorphism
\[
 \phi_B(x)_j=\prod_{i=1}^n x_i^{B_{ij}}
\]
has pullback $\phi_B^*(x^\alpha)=x^{B\alpha}$. Its inverse is the
monomial map associated with $B^{-1}\in\operatorname{GL}_n(\Z)$, so the
jet assertion follows by change of coordinates.
\end{proof}

For compression in the first coordinate, write an exponent as $(a,\beta)$
with $a\in\Z$ and $\beta\in\Z^{n-1}$. For a finite set $S$, define its
row space at $\beta$ by
\[
 W_\beta(S)=\Span_\C\{z^a:(a,\beta)\in S\},
\]
and define the compressed exponent set by
\begin{equation}\label{eq:coordinate-compression}
 \mathcal C_1(S)=
 \{(\ord_{z=1}f,\beta):0\ne f\in W_\beta(S)\}.
\end{equation}
All functions in a row are holomorphic near $z=1$, even when their
exponents are negative. Compression in another coordinate is defined by
permuting coordinates.

\begin{lemma}[Coordinate compression]\label{lem:compression}
Coordinate compression preserves the cardinalities and the graded support
property of a family $(S_k)$. It transfers jet separation backward. In a
nonempty row with smallest and largest exponents $a_-$ and $a_+$, every
compressed exponent lies in $[0,a_+-a_-]$.

Consequently, suppose that, after fixing all but two coordinates, the pair
exponents lie in a scaled convex polygon $\kappa Q$, where $\kappa\geq0$.
Compression in one pair coordinate bounds the new pair exponents by the
shape obtained by moving each interval slice parallel to that coordinate
to start at zero, without changing its length. This shape operation commutes
with scaling by $\kappa$.
\end{lemma}

\begin{proof}
Use the coordinate $u=z-1$ near $z=1$. Lemma~\ref{lem:least-terms}
identifies the possible orders in a row with a basis of distinct leading
orders. If the input sets are graded, the product of two row spaces lies in
the summed degree and remaining exponent. Lemma~\ref{lem:row-valuation}
therefore proves both cardinality preservation and the graded support
property.

Multiplying a nonzero row element by the unit $z^{-a_-}$ gives a polynomial
of degree at most $a_+-a_-$. Its order at $1$ is unchanged and cannot exceed
that degree. This proves the row bound, including for Laurent exponents.

For a fixed degree, choose in every nonempty row a basis with distinct
orders and leading coefficients $1$ in $u$. Substitute $z=1+sZ$, normalize
each basis element by $s$ to its order, and carry along the unchanged
monomial in the other variables. Lemmas~\ref{lem:weighted-limit}
and~\ref{lem:rank-transfer} apply to all rows simultaneously. At a fixed
finite torus tuple, the substitution has nonzero values and is injective
for sufficiently small $s>0$. It gives a local coordinate change, so jet
separation transfers to the original space at the image tuple.

Finally, a row contained in an interval slice of $\kappa Q$ has exponent
span no greater than the length of that slice. Its compressed row is
therefore contained in the interval beginning at zero with that length.
For $\kappa>0$, slices and their lengths scale by $\kappa$. When
$\kappa=0$, the only possible pair exponent is zero and compression leaves
it unchanged.
\end{proof}

When these operations act on a graded family, the support sets are defined
by fixed exact orders and fixed lattice maps. The weights and small
parameters used to transfer jets may nevertheless depend on the degree:
they only establish finite-rank implications. The multiplication property
in Lemma~\ref{lem:row-valuation} holds independently of those analytic
choices.
